% Two consequences of Lemma lem:kernel-l1 (l1_kernel_lemma.tex), written out line by line.
% Fragment in the notation of Appendix G of "From Correlations to Geometry"; it uses the macros and labels
% of the paper (lem:consistency, thm:kernel, prop:reduction, lem:pibound, cor:secondstage,
% lem:kernelmatrix, thm:eigtorus) and does not compile on its own.
% Written on 19 November 2026 by the second session (the repository owner's side).
% STATUS: checked by its author session only; NOT yet read by another session; not formalised; not read by
% a person; not part of the paper. The constants C, C_2', C_7, C_8, C_9 depend only on d.

\begin{proposition}[The reduction in every dimension]\label{prop:reduction-alld}
In the situation and with the notation of Proposition~\ref{prop:reduction}, on the event $E$ and for all
integers $1\le k<(C_1\delta')^{-2}$,
\[
 \big|\Tr P^{k+1}-m_{k+1}-\Pi_k\big|\le C_2'\,\delta'^2\,k\,m_{k+1},
\]
in every dimension $d\ge1$.
\end{proposition}

\begin{proof}
Steps~1, 2, 3(a) and 3(c) of the proof of Proposition~\ref{prop:reduction} are used unchanged. They give
the identity for $\Tr P^{k+1}-m_{k+1}$, the bounds
\begin{equation}\label{eq:alld-hh}
 \big|(h-\hat h)^{(i)}_m(b)\big|\le C\delta'^2m^{-1/2}p_{6m}(x_b-x_i),\qquad
 \big|\hat h^{(i)}_m(b)\big|\le C\delta'm^{-1/2}p_{6m}(x_b-x_i)
\end{equation}
for $1\le m\le k$ (the second from (ii) of Lemma~\ref{lem:consistency}, the first bound and $\delta'\le1$),
and the estimates $C\delta'^2\sqrt k\,m_k$ for the last step of the walk and $C\delta'^2m_k$ for the term
$m=k$. By \eqref{eq:alld-hh}, $p_{6m}\le m_{6m}$ and (iii) of Lemma~\ref{lem:consistency},
\begin{equation}\label{eq:alld-norms}
 \max_{i,b}|h-\hat h|\le C\delta'^2m^{-1/2}m_{6m},\quad\max_i\frac1N\sum_b|h-\hat h|\le C\delta'^2m^{-1/2},\quad
 \max_{i,b}|\hat h|\le C\delta'm^{-1/2}m_{6m},\quad\max_i\frac1N\sum_b|\hat h|\le C\delta'm^{-1/2}.
\end{equation}
It remains to estimate, for $1\le m\le k-1$ and $l=k-m\ge1$, the two sums
\[
 A_m=\frac1N\sum_{i,b}(P^l)_{ib}\,(h-\hat h)^{(i)}_m(b),\qquad
 B_m=\frac1N\sum_{i,b}\Big[(P^l)_{ib}-\frac{p_l(x_i-x_b)}N\Big]\hat h^{(i)}_m(b),
\]
whose sum is the difference in Step~3(b). Recall that $s\mapsto m_s$ is non-increasing and
$m_s\le2^{d/2}m_{2s}$, that $D_l\le C_1\delta'\sqrt l\,m_l\le m_l$ for $l\le k$, and that $l\le k\le\delta'^{-2}$
because $C_1\ge1$.

\emph{The case $l>k/2$.} As in the proof of Proposition~\ref{prop:reduction}, with
$N(P^l)_{ib}\le m_l+D_l\le2m_l$ and $m_l\le m_{k/2}\le2^{d/2}m_k$,
\[
 |A_m|\le2m_l\max_i\frac1N\sum_b|h-\hat h|\le C\delta'^2m^{-1/2}m_k,\qquad
 |B_m|\le D_l\max_i\frac1N\sum_b|\hat h|\le C\delta'^2\sqrt l\,m^{-1/2}m_k .
\]
\emph{The case $l\le k/2$.} Then $m\ge k/2$, so $6m\ge3k$ and $m_{6m}\le m_k$. The rows of $P^l$ are
nonnegative with sum one, hence $|A_m|\le\max_{i,b}|h-\hat h|\le C\delta'^2m^{-1/2}m_k$. For $B_m$ take the
supremum of $\hat h$ out and sum over the first index first:
\[
 |B_m|\le\max_{i,b}\big|\hat h^{(i)}_m(b)\big|\cdot\frac1N\sum_b\Big[\frac1N\sum_i\big|N(P^l)_{ib}-p_l(x_i-x_b)\big|\Big]
 \le C\delta'm^{-1/2}m_k\cdot C_3\,\delta'\sqrt l ,
\]
by Lemma~\ref{lem:kernel-l1}, applied with the number of steps $l$, the fixed second index $b$ and the
sum over the first index $i$ (and $p_l(x_i-x_b)=p_l(x_b-x_i)$).

\emph{Summation.} In both cases $|A_m|+|B_m|\le C\delta'^2\sqrt l\,m^{-1/2}m_k$, and
\[
 \sum_{m=1}^{k-1}\frac{\sqrt{k-m}}{\sqrt m}\le\sqrt k\sum_{m=1}^{k-1}m^{-1/2}\le2k .
\]
Adding the terms of Steps~3(a) and~3(c) and using $m_k\le2^{d/2}m_{k+1}$ gives the claim.
\end{proof}

\begin{corollary}[The second stage in every dimension]\label{cor:secondstage-alld}
Corollary~\ref{cor:secondstage} holds in every dimension $d\ge1$ with $\varphi_d\equiv1$: if $\xi_N\to0$
and $N\xi_N^{d+2}/\log N\to\infty$, then for every sequence $t_N\in[\xi_N^2,1]$
\[
 \frac{\Tr e^{-t_N\Lc}-Ne^{-\tau}}{T_\sigma(t_N)}=1+O_{\mathbb P}\Big(\frac{M\log N}{N}\,J\Big),
\]
and the error is $o_{\mathbb P}(\xi_N^2)$ if $N\xi_N^{d+4}/\log N\to\infty$.
\end{corollary}

\begin{proof}
The proof of Corollary~\ref{cor:secondstage} with Proposition~\ref{prop:reduction-alld} in place of
Proposition~\ref{prop:reduction}; Lemma~\ref{lem:pibound} does not depend on the dimension. For the last
statement: if $\log N\le\xi_N^{-2}$, then $J\le19\xi_N^{-2}$ and the error is $O(\log N/N\xi_N^{d+2})$;
otherwise $J\le20\log N$, $M\le C(\log N)^{d/2}$ and $\xi_N^2>1/\log N$, so the error divided by $\xi_N^2$ is
at most $C(\log N)^{d/2+3}/N$.
\end{proof}

\begin{theorem}[The eigenvalues on a flat torus, with the operator norm from the mean]\label{thm:eigtorus-l1}
In the situation of Theorem~\ref{thm:eigtorus}, on the intersection of the events of
Lemmas~\ref{lem:consistency} and~\ref{lem:kernelmatrix}:
\begin{enumerate}
\item[(i)] for all integers $1\le k\le\delta'^{-2}$ and all $1\le j\le N$,
\[
 \big|s_j^k-e^{-k\sigma\lambda_j}\big|\le C_7\,\delta'\sqrt k+C\sqrt{\frac{\log N}N}\;m_k ;
\]
\item[(ii)] for all $j\le N$ with $\lambda_j\le\xi^{-2}$ and, with $\lambda_*=1\vee\lambda_j$,
\begin{equation}\label{eq:l1-cond}
 C_8\Big[\frac{\delta'}\xi\,\lambda_*^{-1/2}+\sqrt{\frac{\log N}N}\;\lambda_*^{d/2}\Big]\le1,
\end{equation}
\[
 \Big|\lambda^N_j-\frac{1-e^{-\sigma\lambda_j}}{\sigma}\Big|\le C_9\Big[\frac{\delta'}\xi\,\lambda_*^{1/2}
 +\sqrt{\frac{\log N}N}\;\lambda_*^{(d+2)/2}\Big].
\]
\end{enumerate}
\end{theorem}

\begin{proof}
\emph{(i).} $S^k-Q_k$ is real and symmetric, so its operator norm is at most its largest column sum. As in
the proof of Theorem~\ref{thm:eigtorus}, $N(S^k)_{ai}=\rho_{ai}\,N(P^k)_{ai}$ with
$\rho_{ai}=(D_a/D_i)^{1/2}$ and $|\rho_{ai}-1|\le8\delta$, hence
\[
 \sum_a\big|(S^k-Q_k)_{ai}\big|\le8\delta\sum_a(P^k)_{ai}+\frac1N\sum_a\big|N(P^k)_{ai}-p_k(x_a-x_i)\big| .
\]
With the stationary weight $w$ of the proof of Lemma~\ref{lem:kernel-l1}, $\sum_aw_a(P^k)_{ai}=w_i$, so
that $\sum_a(P^k)_{ai}\le w_i/\min_bw_b\le\frac{1+\delta}{1-\delta}\le\frac97$. By
Lemma~\ref{lem:kernel-l1} the second term is at most $C_3\delta'\sqrt k$ for $1\le k\le\delta'^{-2}$. With
$\delta\le6\delta'$,
\[
 \|S^k-Q_k\|\le\tfrac{72}7\,\delta+C_3\,\delta'\sqrt k\le C_7\,\delta'\sqrt k,\qquad C_7=62+C_3 .
\]
$S$ is positive semidefinite, so the eigenvalues of $S^k$ in decreasing order are $s_j^k$; Weyl's inequality
and Lemma~\ref{lem:kernelmatrix} give (i).

\emph{(ii).} Let $a=e^{-\sigma\lambda_j}$ and $k=\lfloor1/\sigma\lambda_*\rfloor$. As in the proof of
Theorem~\ref{thm:eigtorus}: $\sigma\lambda_*\le\frac12$, $k\ge2$, $\frac12\le k\sigma\lambda_*\le1$,
$a^k\ge e^{-1}$, $\sqrt k\le(\sigma\lambda_*)^{-1/2}=\sqrt2\,\xi^{-1}\lambda_*^{-1/2}$ and $m_k\le C\lambda_*^{d/2}$. Hence
the right-hand side of (i) is
\[
 \eta\le\sqrt2\,C_7\,\frac{\delta'}\xi\,\lambda_*^{-1/2}+C\sqrt{\frac{\log N}N}\;\lambda_*^{d/2}.
\]
Two conditions are needed. First $k\le\delta'^{-2}$, the range of (i): since $k\le2/\xi^2\lambda_*$, it holds
if $\sqrt2\,\frac{\delta'}\xi\lambda_*^{-1/2}\le1$. Second $a^k\ge2\eta$, which holds if $2e\,\eta\le1$. Both
follow from \eqref{eq:l1-cond} with $C_8=\max(\sqrt2,\,2\sqrt2\,e\,C_7,\,2eC)$. Then, by the step from
powers to eigenvalues in the proof of Theorem~\ref{thm:eigtorus} (which uses only $a^k\ge2\eta$ and
$|s_j^k-a^k|\le\eta$),
\[
 \Big|\lambda^N_j-\frac{1-a}\sigma\Big|\le\frac{2\,\eta\,e^{k\sigma\lambda_j}}{k\sigma}\le4e\,\lambda_*\,\eta ,
\]
because $k\sigma\ge1/2\lambda_*$ and $k\sigma\lambda_j\le1$. This is (ii) with $C_9=4e\max(\sqrt2\,C_7,C)$.
\end{proof}

\begin{remark}[What changes against Theorem~\ref{thm:eigtorus}]\label{rem:eigtorus-l1}
(i) The random error $C\frac{\delta'}\xi\lambda_*^{(d+1)/2}$ of Theorem~\ref{thm:eigtorus}(ii) is replaced
by $C\frac{\delta'}\xi\lambda_*^{1/2}+C\sqrt{\log N/N}\,\lambda_*^{(d+2)/2}$. The first term is the old bound
times $\lambda_*^{-d/2}$. The second is the old bound times $\xi\lambda_*^{1/2}M^{-1/2}$, which is at most
$C\xi^{d/2}$ in the range $\lambda_*\le\xi^{-2}$, because $M\ge(2\pi\xi^2)^{-d/2}$. The first term comes from the kernel of the
walk (Lemma~\ref{lem:kernel-l1}), the second from the spectrum of the kernel matrix
(Lemma~\ref{lem:kernelmatrix}, a bound in the Hilbert--Schmidt norm, probably not sharp).

(ii) Condition \eqref{eq:l1-cond} holds for all $\lambda_j\le\xi^{-2}$ at once if
$C_8\big[\frac{\delta'}\xi+\sqrt{\log N/N}\,\xi^{-d}\big]\le1$. In the form of
Remark~\ref{rem:eigtorus}(ii): for $1\le\lambda_j\le\xi^{-2}$,
\[
 |\lambda^N_j-\lambda_j|\le\Big(C_9\Big[\frac{\delta'}\xi\,\lambda_j^{-1/2}+\sqrt{\frac{\log N}N}\,\lambda_j^{d/2}\Big]
 +\tfrac14\xi^2\lambda_j\Big)\lambda_j\le\big(\eta_N+\tfrac14\xi^2\lambda_j\big)\lambda_j,\qquad
 \eta_N=C_9\Big[\frac{\delta'}\xi+\sqrt{\frac{\log N}N}\,\xi^{-d}\Big],
\]
and $\eta_N\to0$ if and only if $N\xi_N^{\max(d+2,\,2d)}/\log N\to\infty$ (in place of
$N\xi_N^{3d}/\log N\to\infty$): the same condition in dimension one, $N\xi^4$ against $N\xi^6$ in dimension
two, $N\xi^6$ against $N\xi^9$ in dimension three. Whether this is the form \eqref{eq:eigrate2} of the
paper has not been checked by the second session, which has not seen that equation.

(iii) Corollaries~\ref{cor:eigspace} and~\ref{cor:deloc} use $\|S^k-Q_k\|\le16\delta m_k+D_k$ and the
entries of $S^{2k}$; with (i) the approximate-eigenvector estimate in Corollary~\ref{cor:eigspace} becomes
$\eta_1\le C\delta'm_k+C_7\delta'\sqrt k$. This has not been carried through.

(iv) \emph{Numerical check} (\texttt{scripts/operator\_norm\_check.py}, one sample, $d=2$, $\xi=0.1$,
$N=2500$, $\delta'=0.223$): $\|S^k-Q_k\|=0.091$, $0.061$, $0.046$, $0.031$, $0.023$ for $k=1,2,4,8,16$,
against $\delta'\sqrt k=0.22,\dots,0.89$ and $\delta'\sqrt k\,m_k=3.55,\dots,1.05$. Compatible with (i); in
this run $\delta'$ is not small.
\end{remark}
